% ============================================================================
%  บทสรุปผู้บริหาร (Executive Summary)
% ============================================================================

\chapter*{บทสรุปผู้บริหาร}
\addcontentsline{toc}{chapter}{บทสรุปผู้บริหาร}

\section*{ภาพรวมโครงงาน}

โครงงานนี้มีวัตถุประสงค์เพื่อออกแบบและพัฒนาหุ่นยนต์หยิบและวางชิ้นงานแบบวงกลม
สำหรับลำเลียงชิ้นงาน Connection Rod ระหว่างสี่สถานีทำงานในกระบวนการผลิตเชิงอุตสาหกรรม
ดำเนินการโดยนักศึกษากลุ่มที่ 6 ภายใต้วิชา FRA263/264: Robotics Studio III
สถาบันวิทยาการหุ่นยนต์ภาคสนาม มหาวิทยาลัยเทคโนโลยีพระจอมเกล้าธนบุรี ปีการศึกษา 2568

\section*{สิ่งที่พัฒนาและดำเนินการสำเร็จ}

ทีมพัฒนาระบบได้ดำเนินการออกแบบและสร้างระบบครบทุกองค์ประกอบ ดังนี้

\textbf{ด้านเชิงกล:}
ออกแบบและประกอบแขนกลแบบ 1-DOF ที่ขับเคลื่อนผ่านชุดเฟืองดาวเคราะห์และสายพานฟัน
อัตราทดรวม $N=70$ โดยผ่านการคำนวณตามมาตรฐาน ASME สำหรับเพลา
และ ISO~281 สำหรับตลับลูกปืน พร้อมทั้งวิเคราะห์การโก่งตัวด้วย FEA

\textbf{ด้านไฟฟ้าและอิเล็กทรอนิกส์:}
ออกแบบตู้ควบคุมครบวงจร ประกอบด้วย PCB ที่ออกแบบเอง, วงจรป้องกัน E-Stop,
Optocoupler, Relay และระบบจ่ายไฟ 24~V/10~A, 5~V/5~A พร้อม Circuit Breaker และ Fuse

\textbf{ด้านการควบคุม:}
ระบุพารามิเตอร์ของมอเตอร์ด้วยวิธี System Identification จากข้อมูล Response จริง
ออกแบบตัวควบคุม Cascade PID ด้วย Root Locus บน MATLAB
ติดตั้ง ZV Input Shaper สำหรับลด Residual Vibration ของชิ้นงาน
และออกแบบ Kalman Filter สี่สถานะสำหรับประมาณค่า Disturbance Torque แบบ Real-time

\textbf{ด้านซอฟต์แวร์และการสื่อสาร:}
พัฒนา Firmware บน STM32G474RE ครบทุก module ได้แก่
Cascade PID, ZV Shaper, Kalman Filter, Finite State Machine,
Modbus RTU (LPUART1, Slave Address~21) และ Joystick Interface ผ่าน ESP32

\section*{ข้อจำกัดและปัญหาที่พบ}

\noindent\textbf{1. ปัญหาการเชื่อมต่อกับ Base System (Demo 1)}

\noindent\textbf{Abnormal Heartbeat ใน Demo 1:}
Web Dashboard ของระบบส่ง Heartbeat ไปยัง Robot ไม่สม่ำเสมอ
ส่งผลให้ Robot ตรวจจับ Heartbeat Timeout และตอบกลับ Abnormal Status
ทำให้ Base System แสดงผล Abnormal ตามไปด้วย
สังเกตได้ว่าเมื่อหุ่นยนต์เคลื่อนที่จนค่า Encoder เปลี่ยนแปลง
ปัญหาจะหายชั่วคราวแล้วกลับมาอีก
บ่งชี้ว่าเกิดจาก Timing Mismatch ระหว่างรอบการ Poll ของ Dashboard
กับรอบ Heartbeat Watchdog ของ Firmware
แก้ไขโดยเปลี่ยน Backend เป็น Python asyncio WebSocket
ซึ่งจัดการ Event Loop ได้แม่นยำกว่า
บทเรียน: ควรทดสอบ End-to-End Integration กับ Base System จริง
ก่อน Demo อย่างน้อย 1 สัปดาห์

\medskip
\noindent\textbf{2. ความเสียหายของ Gripper ระหว่างการสาธิต (Demo 2)}

\noindent
ในการสาธิต Demo~2 เกิดปัญหา Abnormal ซ้ำหลังจาก Base System
รัน Process เพื่อสำรวจตำแหน่งวาง Connection Rod เสร็จสิ้น
ทำให้ต้องทำการ Restart Website ส่งผลให้ค่าตำแหน่ง Pick \& Place
ที่ระบบสุ่มมาก่อนหน้าถูก Reset ทั้งหมด
เมื่อพยายามกรอกค่าตำแหน่งใหม่ด้วยตนเอง
พบว่า Base System ไม่ได้แปลงทิศทาง CW/CCW ให้ถูกต้องตามค่าที่กรอก
ส่งผลให้แขนกลเคลื่อนที่ผิดทิศทางและชนกับ Connection Rod
ทำให้ Gripper ได้รับความเสียหายและไม่สามารถซ่อมแซมได้ทัน
จึงไม่สามารถทดสอบ Full Cycle ในการสาธิตรอบสุดท้ายได้
อย่างไรก็ตาม การทำงานของระบบควบคุมและ Firmware ทุก module
ได้รับการยืนยันก่อนเกิดเหตุการณ์ดังกล่าว

\medskip
\noindent\textbf{3. การไม่ได้ติดตั้ง CAN Bus สำหรับ Gripper}

\noindent
ข้อกำหนดเดิมระบุการสื่อสารกับ Gripper ผ่าน CAN Bus
แต่เนื่องจากข้อจำกัดด้านเวลาในการพัฒนา
ทีมจึงตัดสินใจเปลี่ยนมาใช้สัญญาณ Digital I/O โดยตรง
ซึ่งสามารถรองรับการทำงานตามข้อกำหนดสมรรถนะหลักได้ครบถ้วน

\section*{สรุปผลโดยรวม}

ระบบหุ่นยนต์ที่พัฒนาขึ้นครอบคลุมทุกองค์ประกอบของ Cyber-Physical System
ตั้งแต่การออกแบบเชิงกล ไฟฟ้า ซอฟต์แวร์ และการควบคุม
โดยมีการใช้เทคนิคการควบคุมขั้นสูงที่เหมาะสมกับโจทย์อุตสาหกรรม
ความท้าทายหลักที่พบในโครงงานนี้อยู่ที่การ Integration ระหว่างระบบย่อยหลายระบบ
และการบริหารจัดการ Real-time Constraints บน Embedded System
ซึ่งทีมได้แก้ไขและเรียนรู้จากปัญหาดังกล่าวตลอดกระบวนการพัฒนา